%-----------------------------------------------------------------------------
% Załącznik IV
%-----------------------------------------------------------------------------

\cleardoublepage
\section*{IV. Zawadzcy spod Mińska Mazowieckiego}
\label{sec:zawadzcy}

Najstarszym znanym autorowi książki przedstawicielem rodziny Zawadzkich
spod Mińska Mazowieckiego jest Jan
Zawadzki\indeksosoby{Jan}{Zawadzki}{?}{Zakole-Wiktorowo}, urodzony około
1790~roku. Jego żoną była
Apolonia\indeksosoby[Bąk]{Apolonia}{Zawadzka}{?}{Zakole-Wiktorowo} z~domu Bąk,
urodzona około 1795~roku. Nie znamy niestety dokładnych dat ich narodzin czy
ślubu, gdyż akta metrykalne parafii pod wezwaniem Narodzenia Najświętszej
Maryi Panny w~Mińsku Mazowieckim sprzed 1810~roku nie są~publicznie dostępne.
Wiadomo jedynie, iż~do~około 1820~roku mieszkali oni we~wsi
Zakole\footnote{Obecnie wieś ta~nosi nazwę Zakole-Wiktorowo.} znajdującej się
około sześć kilometrów na~południowy wschód od~Mińska Mazowieckiego, gdzie
urodziła się czwórka ich najstarszych dzieci:
\textbf{Józef}\indeksosoby{Józef}{Zawadzki}{1815}{Zakole-Wiktorowo} (ur.~1815~r.~--
zm.~1879~r.), Piotr\indeksosoby{Piotr}{Zawadzki}{1816}{Zakole-Wiktorowo}
(ur.~1816~r.~-- zm.~?),
Tomasz\indeksosoby{Tomasz}{Zawadzki}{1818}{Zakole-Wiktorowo} (ur.~1818~r.~--
zm.~?) oraz
Katarzyna\indeksosoby[Zawadzka][Puławska]{Katarzyna}{Kiliszczyk}{1819}{Zakole-Wiktorowo}
(ur.~1819~r.~-- zm.~1864~r.)~-- patrz: ryc. \ref{fig:jzawadzki_1815},
\ref{fig:pzawadzki_1816}, \ref{fig:tzawadzki_1818} oraz
\ref{fig:kzawadzka_1819}. Po~czym około 1825~roku
Jan\indeksosoby{Jan}{Zawadzki}{?}{Zakole-Wiktorowo}
i~Apolonia\indeksosoby[Bąk]{Apolonia}{Zawadzka}{?}{Zakole-Wiktorowo} Zawadzcy
przeprowadzili się do~nowo utworzonej wsi Anielina\footnote{Co~ciekawe
w~1822~roku z~Zakola do~Anieliny przeprowadził się wraz z~najbliższą rodziną
oraz żoną Apolonią również inny Jan Zawadzki, o~czym wspominałem na~początku
rozdziału czwartego. Jednak jego żona przed ślubem nosiła nazwisko Kossakowska
a~nie Bąk. Nie wiadomo czy obaj imiennicy Zawadzcy byli ze~sobą spokrewnieni,
na~pewno mieszkali razem we~wsi Zakole a~potem w~podobnym czasie
przeprowadzili się do~Anieliny, co~mogłoby wskazywać na~ich bliską relację.},
gdzie na~świat przyszła ich najmłodsza córka
Marianna\indeksosoby{Marianna}{Zawadzka}{1826}{Anielina} (ur.~1826~r.~--
zm.~?)~-- patrz: ryc. \ref{fig:mzawadzka_1826}.

\begin{figure}[!ht]
    \vspace*{0.5cm}
    \centering \includegraphics[width=0.78\linewidth]{
        1815_Józef_Zawadzki_akt_chrztu_parafia_Mińsk_Mazowiecki_wpis_41.jpg}
    \captionsetup{format=hang}
    \caption{Akt chrztu Józefa Zawadzkiego -- par. Mińsk Mazowiecki 1815~rok
    (41/1815) \cite{par_minsk1}.}
    \label{fig:jzawadzki_1815}
\end{figure}

\begin{figure}[!ht]
    \vspace*{0.5cm}
    \centering \includegraphics[width=0.82\linewidth]{
        1816_Piotr_Zawadzki_akt_chrztu_parafia_Mińsk_Mazowiecki_wpis_91.jpg}
    \captionsetup{format=hang}
    \caption{Akt chrztu Piotra Zawadzkiego -- par. Mińsk Mazowiecki 1816~rok
    (91/1816) \cite{par_minsk1}.}
    \label{fig:pzawadzki_1816}
\end{figure}

\begin{figure}[!ht]
    \vspace*{0.5cm}
    \centering \includegraphics[width=0.82\linewidth]{
        1818_Tomasz_Zawadzki_akt_chrztu_parafia_Mińsk_Mazowiecki_wpis_87.jpg}
    \captionsetup{format=hang}
    \caption{Akt chrztu Tomasza Zawadzkiego -- par. Mińsk Mazowiecki 1818~rok
    (87/1818) \cite{par_minsk1}.}
    \label{fig:tzawadzki_1818}
\end{figure}

\begin{figure}[!ht]
    \vspace*{0.5cm}
    \centering \includegraphics[width=0.8\linewidth]{
        1819_Katarzyna_Zawadzka_akt_chrztu_parafia_Mińsk_Mazowiecki_wpis_232.jpg}
    \captionsetup{format=hang}
    \caption{Akt chrztu Katarzyny Zawadzkiej -- par. Mińsk Mazowiecki 1819~rok
    (232/1819) \cite{par_minsk1}.}
    \label{fig:kzawadzka_1819}
\end{figure}

\begin{figure}[!ht]
    \vspace*{0.5cm}
    \centering \includegraphics[width=0.8\linewidth]{
        1826_Marianna_Zawadzka_akt_chrztu_parafia_Mińsk_Mazowiecki_wpis_259.jpg}
    \captionsetup{format=hang}
    \caption{Akt chrztu Marianny Zawadzkiej -- par. Mińsk Mazowiecki 1826~rok
    (259/1826) \cite{par_minsk2}.}
    \label{fig:mzawadzka_1826}
\end{figure}

Jan Zawadzki\indeksosoby{Jan}{Zawadzki}{?}{Zakole-Wiktorowo} zmarł 7~marca
1831~roku w~wieku około czterdziestu lat~-- zaledwie kilka lat
po~przeprowadzce do~Anieliny (patrz: ryc. \ref{fig:jzawadzki_1831}). Jego żona
Apolonia\indeksosoby[Bąk]{Apolonia}{Zawadzka}{?}{Zakole-Wiktorowo} żyła
jeszcze blisko dwadzieścia pięć lat umierając dopiero 23~lutego 1855~roku
w~wieku około sześćdziesięciu lat (patrz: ryc. \ref{fig:azawadzka_1855}).
Oboje do~końca swoich dni mieszkali w~Anielinie.

\begin{figure}[!ht]
    \vspace*{0.5cm}
    \centering \includegraphics[width=0.97\linewidth]{
        1831_Jan_Zawadzki_akt_zgonu_parafia_Mińsk_Mazowiecki_wpis_324.jpg}
    \captionsetup{format=hang}
    \caption{Akt zgonu Jana Zawadzkiego -- par. Mińsk Mazowiecki 1831~rok
    (324/1831) \cite{par_minsk2}.}
    \label{fig:jzawadzki_1831}
\end{figure}

\begin{figure}[!ht]
    \vspace*{0.5cm}
    \centering \includegraphics[width=0.97\linewidth]{
        1855_Apolonia_Zawadzka_akt_zgonu_parafia_Mińsk_Mazowiecki_wpis_84.jpg}
    \captionsetup{format=hang}
    \caption{Akt zgonu Apolonii Zawadzkiej -- par. Mińsk Mazowiecki 1855~rok
    (84/1855) \cite{par_minsk2}.}
    \label{fig:azawadzka_1855}
\end{figure}

Spośród piątki dzieci Jana\indeksosoby{Jan}{Zawadzki}{?}{Zakole-Wiktorowo}
i~Apolonii\indeksosoby[Bąk]{Apolonia}{Zawadzka}{?}{Zakole-Wiktorowo}
Zawadzkich prawdopodobnie tylko dwójka dożyła wieku dorosłego:
Józef\indeksosoby{Józef}{Zawadzki}{1815}{Zakole-Wiktorowo} oraz
Katarzyna\indeksosoby[Zawadzka][Puławska]{Katarzyna}{Kiliszczyk}{1819}{Zakole-Wiktorowo}.
Józef praktycznie całe swoje życie spędził w~rodzinnej Anielinie. Katarzyna
natomiast od~razu po~ślubie wyprowadziła się do~sąsiedniej wsi~Osiny,
potem przeprowadziła się do~Barczącej, gdzie dożyła już reszty swoich dni.
W~niniejszym załączniku skupimy się na~dalszej historii Józefa, gdyż to~jego
dzieci Marcin oraz Marianna były współzałożycielami spółki \emph{Desna} oraz
związały swoje losy z~rodziną Gomulskich z~Desna.

\begin{figure}[!ht]
    \vspace*{0.5cm}
    \centering \includegraphics[width=0.97\linewidth]{
        1837_Józef_Zawadzki_Konstancja_Okoła_akt_ślubu_parafia_Mińsk_Mazowiecki_wpis_26.jpg}
    \captionsetup{format=hang}
    \caption{Akt ślubu Józefa Zawadzkiego i~Konstancji Okoły -- par. Mińsk
    Mazowiecki 1837~rok (26/1837) \cite{par_minsk2}.}
    \label{fig:jzawadzki_1837}
\end{figure}

\textbf{Józef Zawadzki}\indeksosoby{Józef}{Zawadzki}{1815}{Zakole-Wiktorowo}
dnia 5~lutego 1837~roku wziął ślub w~parafii pod wezwaniem Narodzenia
Najświętszej Maryi Panny w~Mińsku Mazowieckim z~dwudziestoletnią wówczas
Konstancją Okołą\indeksosoby[Okoła]{Konstacja}{Zawadzka}{1817}{Stara Niedziałka}
(ur.~1817~r.~-- zm.~1855~r.), panną z~sąsiedniej wsi Niedziałka\footnote{Obecnie
wieś ta~nazywa się Stara Niedziałka.}, córką Stefana i~Agaty z~domu
Olszewskiej (patrz: ryc. \ref{fig:jzawadzki_1837}). Para doczekała się
siódemki dzieci:

\begin{itemize}
\item Franciszki\indeksosoby[Zawadzka]{Franciszka}{Kwiatkowska}{1837}{Jarząbek}
(ur.~1837~r.~-- zm.~1883~r.),
\item Marianny\indeksosoby{Marianna}{Zawadzka}{1840}{Anielina} (ur.~1840~r.~--
zm.~1841~r.),
\item Jana\indeksosoby{Jan}{Zawadzki}{1842}{Anielina} (ur.~1840~r.~--
zm.~1843~r.),
\item Jana\indeksosoby{Jan}{Zawadzki}{1844}{Anielina} (ur.~1844~r.~--
zm.~1849~r.),
\item Katarzyny\indeksosoby{Katarzyna}{Zawadzka}{1846}{Anielina}
(ur.~1846~r.~-- zm.~1860~r.),
\item Józefy\indeksosoby{Józefa}{Zawadzka}{1849}{Anielina} (ur.~1849~r.~--
zm.~1855~r.),
\item Józefa\indeksosoby{Józef}{Zawadzki}{1852}{Anielina} (ur.~1852~r.~--
zm.~?).
\end{itemize}

Dnia 16~lutego 1855~roku, po~osiemnastu latach małżeństwa, Konstancja
Zawadzka\indeksosoby[Okoła]{Konstacja}{Zawadzka}{1817}{Stara Niedziałka}
zmarła (patrz: ryc. \ref{fig:kzawadzka_1855}), pozostawiając po~sobie
owdowiałego męża Józefa\indeksosoby{Józef}{Zawadzki}{1815}{Zakole-Wiktorowo}
oraz czwórkę niepełnoletnich dzieci: Franciszkę, Katarzynę, Józefę oraz Józefa.

\begin{figure}[!ht]
    \vspace*{0.5cm}
    \centering \includegraphics[width=0.92\linewidth]{
        1855_Konstancja_Zawadzka_akt_zgonu_parafia_Mińsk_Mazowiecki_wpis_75.jpg}
    \captionsetup{format=hang}
    \caption{Akt zgonu Konstancji Zawadzkiej -- par. Mińsk Mazowiecki 1855~rok
    (75/1855) \cite{par_minsk2}.}
    \label{fig:kzawadzka_1855}
\end{figure}

\begin{figure}[!ht]
    \vspace*{0.5cm}
    \centering \includegraphics[width=0.92\linewidth]{
        1855_Józef_Zawadzki_Marianna_Górka_akt_ślubu_parafia_Mińsk_Mazowiecki_wpis_30.jpg}
    \captionsetup{format=hang}
    \caption{Akt ślubu Józefa Zawadzkiego i~Marianny Górki -- par. Mińsk
    Mazowiecki 1855~rok (30/1855) \cite{par_minsk2}.}
    \label{fig:jzawadzki_1855}
\end{figure}

Józef Zawadzki\indeksosoby{Józef}{Zawadzki}{1815}{Zakole-Wiktorowo} wziął
ponownie ślub zaledwie siedem miesięcy po~śmierci swojej pierwszej żony~--
9~września 1855~roku w~wieku czterdziestu lat ożenił się
z~dwudziestoośmioletnią Marianną
Górką\indeksosoby[Górka]{Marianna}{Zawadzka}{1827}{Rudno} (ur.~1827~r.~--
zm.~?), panną z~Rudna, córką Wincentego i~Salomei z~domu Wyszogrodzkiej
(patrz: ryc. \ref{fig:jzawadzki_1855}). Nowi małżonkowie Zawadzcy doczekali
się prawdopodobnie szóstki dzieci:

\begin{itemize}
\item \textbf{Marcina}\indeksosoby{Marcin}{Zawadzki}{1856}{Anielina}
(ur.~1856~r.~-- zm.~1927~r.),
\item Jana\indeksosoby{Jan}{Zawadzki}{1859}{Anielina} (ur.~1859~r.~-- zm.~?),
\item Rozalii\indeksosoby{Rozalia}{Zawadzka}{1863}{Anielina} (ur.~1863~r.~--
zm.~1879~r.),
\item Antoniego\indeksosoby{Antoni}{Zawadzki}{1865}{Anielina} (ur.~1865~r.~--
zm.~?),
\item Marianny\indeksosoby[Gruba]{Marianna}{Zawadzka}{1867}{Anielina}
(ur.~1867~r.~-- zm.~1943~r.),
\item Józefa\indeksosoby{Józef}{Zawadzki}{1872}{Anielina} (ur.~1872~r.~--
zm.~1903~r.).
\end{itemize}

Dnia 21~grudnia 1879~roku w~wieku 64~lat Józef 
Zawadzki\indeksosoby{Józef}{Zawadzki}{1815}{Zakole-Wiktorowo} zmarł~--
patrz: ryc. \ref{fig:jzawadzki_1879}. Data zgonu jego drugiej żony Marianny
nie jest niestety znana autorowi książki.

\begin{figure}[!ht]
    \vspace*{0.5cm}
    \centering \includegraphics[width=0.97\linewidth]{
        1856_Marcin_Zawadzki_akt_chrztu_parafia_Mińsk_Mazowiecki_wpis_144.jpg}
    \captionsetup{format=hang}
    \caption{Akt chrztu Marcina Zawadzkiego -- par. Mińsk Mazowiecki 1856~rok
    (144/1856) \cite{par_minsk2}.}
    \label{fig:mzawadzki_1856}
\end{figure}

\begin{figure}[!ht]
    \vspace*{0.5cm}
    \centering \includegraphics[width=0.97\linewidth]{
        1867_Marianna_Zawadzka_akt_chrztu_parafia_Mińsk_Mazowiecki_wpis_239.jpg}
    \captionsetup{format=hang}
    \caption{Akt chrztu Marianny Zawadzkiej -- par. Mińsk Mazowiecki 1867~rok
    (239/1867) \cite{par_minsk2}.}
    \label{fig:mzawadzki_1867}
\end{figure}

\begin{figure}[!ht]
    \vspace*{0.5cm}
    \centering \includegraphics[width=0.97\linewidth]{
        1879_Józef_Zawadzki_akt_zgonu_parafia_Długa_Kościelna_wpis_219.jpg}
    \captionsetup{format=hang}
    \caption{Akt zgonu Józefa Zawadzkiego -- par. Mińsk Mazowiecki 1879~rok
    (219/1879) \cite{par_minsk2}.}
    \label{fig:jzawadzki_1879}
\end{figure}

Jak już wspominano wcześniej, spośród trzynaściorga dzieci Józefa
Zawadzkiego\indeksosoby{Józef}{Zawadzki}{1815}{Zakole-Wiktorowo} dwójka
osiedliła się około 1894~roku w~nowo powstałej wsi \emph{Desna}~-- 
Marcin\indeksosoby{Marcin}{Zawadzki}{1856}{Anielina} oraz jego młodsza siostra
Marianna\indeksosoby[Gruba]{Marianna}{Zawadzka}{1867}{Anielina}. W~przypadku
Marcina była to~jego docelowa miejscowość zamieszkania. Natomiast Marianna
wraz z~mężem Marcinem Grubą\indeksosoby{Marcin}{Gruba}{1858}{Wólka Barcząca}
po~kilku latach mieszkania na~\emph{Deśnie} sprzedali swoją ziemię Michałowi
Gomulskiemu\indeksosoby{Michał}{Gomulski}{1859}{Anielina} i~przeprowadzili
z~powrotem w~okolice Mińska Mazowieckiego.

\begin{figure}[!ht]
    \vspace*{0.5cm}
    \centering \includegraphics[width=1.0\linewidth]{
        1881_Marcin_Zawadzki_Marianna_Zawadzka_akt_ślubu_parafia_Mińsk_Mazowiecki_wpis_18.jpg}
    \captionsetup{format=hang}
    \caption{Akt ślubu Marcina Zawadzkiego i~Marianny Zawadzkiej -- par. Mińsk
    Mazowiecki 1881~rok (18/1881) \cite{par_minsk2}.}
    \label{fig:mzawadzki_1881}
\end{figure}

Dnia 16~lutego 1881~roku
\textbf{Marcin Zawadzki}\indeksosoby{Marcin}{Zawadzki}{1856}{Anielina} ożenił
się z~szesnastoletnią wówczas Marianną
Zawadzką\indeksosoby[Zawadzka]{Marianna}{Zawadzka}{1864}{Targówka}
(ur.~1864~r.~-- zm.~1939~r.), panną ze~wsi Targówka, córką Jana i~Anny z~domu
Karpińskiej (patrz: ryc. \ref{fig:mzawadzki_1881}). Pomimo iż~nowożeńcy
Zawadzcy nosili przed ślubem to~samo nazwisko, a~nawet wychowywali się
w~sąsiednich wsiach pod Mińskiem Mazowieckim, to~nie byli ze~sobą w~żaden
sposób spokrewnieni. Jednak zbieżność ich nazwisk nie była też całkowicie
przypadkowa.

\begin{figure}[!ht]
    \vspace*{0.5cm}
    \centering \includegraphics[width=0.97\linewidth]{
        1839_Jan_Zawadka_akt_chrztu_parafia_Mińsk_Mazowiecki_wpis_60.jpg}
    \captionsetup{format=hang}
    \caption{Akt chrztu Jana Zawadki -- par. Mińsk Mazowiecki 1839~rok
    (60/1839) \cite{par_minsk2}.}
    \label{fig:jzawadzki_1839}
\end{figure}

Pradziadek Marianny\indeksosoby[Zawadzka]{Marianna}{Zawadzka}{1864}{Targówka}
Wojciech\indeksosoby{Wojciech}{Zawadka}{1776}{Starogród} około 1825~roku
przeprowadził się wraz z~rodziną ze~wsi Starogród pod Siennicą
do~nowo~powstałej wsi~Anielina\footnote{W~tym samym czasie do~Anieliny
z~Zakola sprowadził się dziadek Marcina~-- Jan, co~szerzej opisano na~początku
bieżącego załącznika.}. Nosił on~wówczas nazwisko Zawadka, podobnie wszystkie
jego dzieci były chrzczone pod tym nazwiskiem, a~także wnuki, w~tym co~istotne
ojciec Marianny Jan\indeksosoby{Jan}{Zawadka}{1839}{Anielina}~-- patrz: ryc.
\ref{fig:jzawadzki_1839}. Dopiero około 1850~roku nazwisko Zawadka zaczęło
stopniowo zanikać z~ksiąg metrykalnych parafii pod wezwaniem Narodzenia
Najświętszej Maryi Panny w~Mińsku Mazowieckim, a~wszyscy ochrzczeni Zawadkowie
zaczęli być określani jako Zawadzcy~-- w~akcie ślubu rodziców Marianny
pojawia się już to~nazwisko (patrz: ryc. \ref{fig:jzawadzki_1858}). Skąd więc
to~zanikanie jednego nazwiska na~rzecz drugiego zbliżonego? Prawdopodobnie
z~faktu, iż~w~połowie XIX~wieku do~mińskiej parafii należało co~najmniej kilka
rodzin Zawadzkich żyjących w~okolicach Mińska Mazowieckiego i~tylko jedna
rodzina Zawadków (ta~wywodząca się od~Wojciecha). Tamtejszy proboszcz
najwyraźniej uznał, iż~poprawny zapis nazwiska to~ten, którym posługiwała się
większa liczba parafian i~zaczął go~konsekwentnie stosować we~wszystkich
sporządzanych przez siebie aktach metrykalnych.

\begin{figure}[!ht]
    \vspace*{0.5cm}
    \centering \includegraphics[width=0.9\linewidth]{
        1858_Jan_Zawadzki_Anna_Karpińska_akt_ślubu_parafia_Mińsk_Mazowiecki_wpis_85.jpg}
    \captionsetup{format=hang}
    \caption{Akt ślubu Jana Zawadzkiego (Zawadki) i~Anny Karpińskiej~-- par.
    Mińsk Mazowiecki 1858~rok (85/1858) \cite{par_minsk2}.}
    \label{fig:jzawadzki_1858}
\end{figure}

\begin{figure}[!ht]
    \vspace*{0.5cm}
    \centering \includegraphics[width=0.8\linewidth]{
        1864_Marianna_Zawadzka_akt_chrztu_parafia_Mińsk_Mazowiecki_wpis_256.jpg}
    \captionsetup{format=hang}
    \caption{Akt chrztu Marianny Zawadzkiej~-- par. Mińsk Mazowiecki 1864~rok
    (256/1864) \cite{par_minsk2}.}
    \label{fig:mzawadzka_1864}
\end{figure}

Marcin\indeksosoby{Marcin}{Zawadzki}{1856}{Anielina}
i~Marianna\indeksosoby[Zawadzka]{Marianna}{Zawadzka}{1864}{Targówka} Zawadzcy
po~ślubie zamieszkali w~rodzinnej wsi pana młodego~-- w~Anielinie, gdzie
urodziła się szóstka ich najstarszych dzieci. Po~około czternastu latach
rodzina Zawadzkich przeprowadziła się do~nowo~powstałej wsi \emph{Desna},
gdzie na~świat przyszła pozostała ósemka ich potomstwa. Łącznie małżonkowie
Zawadzcy dochowali się czternaściorga dzieci (trzech córek i~jedenastu synów): 

\begin{itemize}
\item Józefa\indeksosoby{Józef}{Zawadzki}{1882}{Anielina}
(ur.~1882~r.~-- zm.~?)~-- patrz: ryc. \ref{fig:jzawadzki_1882},
\item Stanisławy\indeksosoby[Zawadzka]{Stanisława}{Wróbel}{1885}{Anielina}
(ur.~1885~r.~-- zm.~?)~-- patrz: ryc. \ref{fig:szawadzka_1885},
\item Wincentego\indeksosoby{Wincenty}{Zawadzki}{1887}{Anielina}
(ur.~1887~r.~-- zm.~?)~-- patrz: ryc. \ref{fig:wzawadzki_1887},
\item Jana\indeksosoby{Jan}{Zawadzki}{1889}{Anielina}
(ur.~1889~r.~-- zm.~?)~-- patrz: ryc. \ref{fig:jzawadzki_1889},
\item Pawła\indeksosoby{Paweł}{Zawadzki}{1891}{Anielina}
(ur.~1891~r.~-- zm.~?)~-- patrz: ryc. \ref{fig:pzawadzki_1891},
\item Michała\indeksosoby{Michał}{Zawadzki}{1893}{Anielina}
(ur.~1893~r.~-- zm.~?)~-- patrz: ryc. \ref{fig:mzawadzki_1893},
\item \textbf{Magdaleny}\indeksosoby[Zawadzka][Gomulska]{Magdalena}{Poziemska}{1895}{Cisie}
(ur.~1895~r.~-- zm.~1970~r.)~-- patrz: ryc. \ref{fig:mzawadzka_1896_2},
\item Andrzeja\indeksosoby{Andrzej}{Zawadzki}{1897}{Desno}
(ur.~1897~r.~-- zm.~1898~r.)~-- patrz: ryc. \ref{fig:azawadzki_1897},
\item Antoniego\indeksosoby{Antoni}{Zawadzki}{1899}{Desno}
(ur.~1899~r.~-- zm.~?)~-- patrz: ryc. \ref{fig:azawadzki_1899},
\item Feliksy\indeksosoby{Feliks}{Zawadzki}{1900}{Desno}
(ur.~1900~r.~-- zm.~?)~-- patrz: ryc. \ref{fig:fzawadzki_1900},
\item Ludwika\indeksosoby{Ludwik}{Zawadzki}{1902}{Desno}
(ur.~1902~r.~-- zm.~1904~r.)~-- patrz: ryc. \ref{fig:lzawadzki_1902},
\item Mariana\indeksosoby{Marian}{Zawadzki}{1904}{Desno}
(ur.~1904~r.~-- zm.~1905~r.)~-- patrz: ryc. \ref{fig:mzawadzki_1904},
\item Juliana\indeksosoby{Julian}{Zawadzki}{1905}{Desno}
(ur.~1905~r.~-- zm.~?)~-- patrz: ryc. \ref{fig:jzawadzki_1905},
\item Cecylii\indeksosoby{Cecylia}{Zawadzka}{1908}{Desno}
(ur.~1908~r.~-- zm.~?)~-- patrz: ryc. \ref{fig:czawadzka_1908}.
\end{itemize}

\begin{figure}[!ht]
    \vspace*{0.5cm}
    \centering \includegraphics[width=0.78\linewidth]{
        1882_Józef_Zawadzki_akt_chrztu_parafia_Mińsk_Mazowiecki_wpis_276.jpg}
    \captionsetup{format=hang}
    \caption{Akt chrztu Józefa Zawadzkiego~-- par. Mińsk Mazowiecki 1882~rok
    (276/1882) \cite{par_minsk2}.}
    \label{fig:jzawadzki_1882}
\end{figure}

\begin{figure}[!ht]
    \vspace*{0.5cm}
    \centering \includegraphics[width=0.78\linewidth]{
        1885_Stanisława_Zawadzka_akt_chrztu_parafia_Mińsk_Mazowiecki_wpis_135.jpg}
    \captionsetup{format=hang}
    \caption{Akt chrztu Stanisławy Zawadzkiej~-- par. Mińsk Mazowiecki
    1885~rok (135/1885) \cite{par_minsk2}.}
    \label{fig:szawadzka_1885}
\end{figure}

\begin{figure}[!ht]
    \vspace*{0.5cm}
    \centering \includegraphics[width=0.76\linewidth]{
        1887_Wincenty_Zawadzki_akt_chrztu_parafia_Mińsk_Mazowiecki_wpis_225.jpg}
    \captionsetup{format=hang}
    \caption{Akt chrztu Wincentego Zawadzkiego~-- par. Mińsk Mazowiecki
    1887~rok (225/1887) \cite{par_minsk2}.}
    \label{fig:wzawadzki_1887}
\end{figure}

\begin{figure}[!ht]
    \vspace*{0.5cm}
    \centering \includegraphics[width=0.75\linewidth]{
        1889_Jan_Zawadzki_akt_chrztu_parafia_Mińsk_Mazowiecki_wpis_86.jpg}
    \captionsetup{format=hang}
    \caption{Akt chrztu Jana Zawadzkiego~-- par. Mińsk Mazowiecki 1889~rok
    (86/1889) \cite{par_minsk2}.}
    \label{fig:jzawadzki_1889}
\end{figure}

\begin{figure}[!ht]
    \vspace*{0.5cm}
    \centering \includegraphics[width=0.75\linewidth]{
        1891_Paweł_Zawadzki_akt_chrztu_parafia_Mińsk_Mazowiecki_wpis_132.jpg}
    \captionsetup{format=hang}
    \caption{Akt chrztu Pawła Zawadzkiego~-- par. Mińsk Mazowiecki 1891~rok
    (132/1891) \cite{par_minsk2}.}
    \label{fig:pzawadzki_1891}
\end{figure}

\begin{figure}[!ht]
    \vspace*{0.5cm}
    \centering \includegraphics[width=0.75\linewidth]{
        1893_Michał_Zawadzki_akt_chrztu_parafia_Mińsk_Mazowiecki_wpis_314.jpg}
    \captionsetup{format=hang}
    \caption{Akt chrztu Michała Zawadzkiego~-- par. Mińsk Mazowiecki 1893~rok
    (314/1893) \cite{par_minsk2}.}
    \label{fig:mzawadzki_1893}
\end{figure}

\begin{figure}[!ht]
    \centering \includegraphics[width=0.85\linewidth]{
        1895_Magdalena_Gomulska_Zawadzka_Poziemska_akt_chrztu_parafia_Długa_Kościelna_wpis_66.jpg}
    \captionsetup{format=hang}
    \caption{Akt chrztu Magdaleny Zawadzkiej -- par. Długa Kościelna 1895~rok
    (66/1895) \cite{par_dluga}.}
    \label{fig:mzawadzka_1896_2}
\end{figure}

\begin{figure}[!ht]
    \centering \includegraphics[width=0.85\linewidth]{
        1897_Andrzej_Zawadzki_akt_chrztu_parafia_Długa_Kościelna_wpis_121.jpg}
    \captionsetup{format=hang}
    \caption{Akt chrztu Andrzeja Zawadzkiego -- par. Długa Kościelna 1897~rok
    (121/1897) \cite{par_dluga}.}
    \label{fig:azawadzki_1897}
\end{figure}

\begin{figure}[!ht]
    \centering \includegraphics[width=0.74\linewidth]{
        1899_Antoni_Zawadzki_akt_chrztu_parafia_Długa_Kościelna_wpis_50.jpg}
    \captionsetup{format=hang}
    \caption{Akt chrztu Antoniego Zawadzkiego -- par. Długa Kościelna 1899~rok
    (50/1899) \cite{par_dluga}.}
    \label{fig:azawadzki_1899}
\end{figure}

\begin{figure}[!ht]
    \centering \includegraphics[width=0.74\linewidth]{
        1900_Feliks_Zawadzki_akt_chrztu_parafia_Długa_Kościelna_wpis_94.jpg}
    \captionsetup{format=hang}
    \caption{Akt chrztu Feliksa Zawadzkiego -- par. Długa Kościelna 1900~rok
    (94/1900) \cite{par_dluga}.}
    \label{fig:fzawadzki_1900}
\end{figure}

\begin{figure}[!ht]
    \centering \includegraphics[width=0.82\linewidth]{
        1902_Ludwik_Zawadzki_akt_chrztu_parafia_Długa_Kościelna_wpis_41.jpg}
    \captionsetup{format=hang}
    \caption{Akt chrztu Ludwika Zawadzkiego -- par. Długa Kościelna 1902~rok
    (41/1902) \cite{par_dluga}.}
    \label{fig:lzawadzki_1902}
\end{figure}

\begin{figure}[!ht]
    \centering \includegraphics[width=0.82\linewidth]{
        1904_Marian_Zawadzki_akt_chrztu_parafia_Długa_Kościelna_wpis_84.jpg}
    \captionsetup{format=hang}
    \caption{Akt chrztu Mariana Zawadzkiego -- par. Długa Kościelna 1904~rok
    (84/1904) \cite{par_dluga}.}
    \label{fig:mzawadzki_1904}
\end{figure}

\begin{figure}[!ht]
    \centering \includegraphics[width=0.85\linewidth]{
        1905_Julian_Zawadzki_akt_chrztu_parafia_Długa_Kościelna_wpis_128.jpg}
    \captionsetup{format=hang}
    \caption{Akt chrztu Juliana Zawadzkiego -- par. Długa Kościelna 1905~rok
    (128/1905) \cite{par_dluga}.}
    \label{fig:jzawadzki_1905}
\end{figure}

\begin{figure}[!ht]
    \centering \includegraphics[width=0.85\linewidth]{
        1908_Cecylia_Zawadzka_akt_chrztu_parafia_Długa_Kościelna_wpis_111.jpg}
    \captionsetup{format=hang}
    \caption{Akt chrztu Cecylii Zawadzkiej -- par. Długa Kościelna 1908~rok
    (111/1908) \cite{par_dluga}.}
    \label{fig:czawadzka_1908}
\end{figure}

Marcin Zawadzki\indeksosoby{Marcin}{Zawadzki}{1856}{Anielina} zmarł 29~marca
1927~roku w~wieku 71~lat pozostawiając po~sobie owdowiałą żonę oraz dziewiątkę
dorosłych już dzieci\footnote{Spośród czternaściorga dzieci Marcina i~Marianny
Zawadzkich prawdopodobnie tylko dziewiątka dożyła wieku dorosłego.}.
Marianna\indeksosoby[Zawadzka]{Marianna}{Zawadzka}{1864}{Targówka} Zawadzka
zmarła 9~października 1939~roku w~wieku 75~lat. Oboje zostali pochowani
na~cmentarzu parafialnym w~Długiej Kościelnej~-- patrz ryc.
\ref{fig:mzawadzki_1927} oraz \ref{fig:mzawadzka_1939}.

\begin{figure}[!ht]
    \vspace*{0.4cm}
    \centering \includegraphics[width=1.0\linewidth]{
       1927_Marcin_Zawadzki_zdjęcie_nagrobka1_cmentarz_Długa_Kościelna.jpg}
    \captionsetup{format=hang}
    \caption{Zdjęcie grobu Marcina Zawadzkiego, znajdującego się na~cmentarzu
    parafii św. Anny w~Długiej Kościelnej (sektor: A5, rząd: 9, numer: 13).}
    \label{fig:mzawadzki_1927}
\end{figure}

\begin{figure}[!ht]
    \vspace*{0.4cm}
    \centering \includegraphics[width=1.0\linewidth]{
       1939_Marianna_Zawadzka_zdjęcie_nagrobka2_cmentarz_Długa_Kościelna.jpg}
    \captionsetup{format=hang}
    \caption{Zdjęcie grobu Marianny Zawadzkiej, znajdującego się na~cmentarzu
    parafii św.~Anny w~Długiej Kościelnej (sektor: A5, rząd: 10, numer: 15).}
    \label{fig:mzawadzka_1939}
\end{figure}

Po~śmierci małżonków Zawadzkich cały ich majątek na~Deśnie (około 8~hektarów
ziemi) został podzielony na~dziewięć równych części pomiędzy:
Józefa\indeksosoby{Józef}{Zawadzki}{1882}{Anielina},
Stanisławę\indeksosoby[Zawadzka]{Stanisława}{Wróbel}{1885}{Anielina},
Wincentego\indeksosoby{Wincenty}{Zawadzki}{1887}{Anielina},
Jana\indeksosoby{Jan}{Zawadzki}{1889}{Anielina},
Pawła\indeksosoby{Paweł}{Zawadzki}{1891}{Anielina},
Magdalenę\indeksosoby[Zawadzka][Gomulska]{Magdalena}{Poziemska}{1895}{Cisie}
Antoniego\indeksosoby{Antoni}{Zawadzki}{1899}{Desno},
Feliksa\indeksosoby{Feliks}{Zawadzki}{1900}{Desno},
oraz Cecylię\indeksosoby{Cecylia}{Zawadzka}{1908}{Desno}. Informuje o~tym wpis
z~1944~roku znajdujący się w~księdze hipotecznej wsi \emph{Desna}
\enquote{Dobra ziemskie DESNA i~KĄTY MROWISKIE t.~1, powiat nowomiński,
wł.~Izrael Rotgolc} \cite{desna1} należącej do~zbioru \enquote{Hipoteka
w~Mińsku Mazowieckim} Archiwum Państwowego w~Siedlcach (patrz: ryc.
\ref{fig:mzawadzki_1944_1}).

\begin{figure}[!ht]
    \centering
    \includegraphics[width=0.9\linewidth]{Zawadzcy_podział_ziemi_1.JPG}
    \captionsetup{format=hang}
    \caption{Fragment wpisu w~księdze hipotecznej wsi \emph{Desna} informujący
    o~ostatecznym podziale majątku Marcina i~Marianny Zawadzkich - część 1.
    \cite{desna1}.}
    \label{fig:mzawadzki_1944_1}
\end{figure}

\begin{figure}[!ht]
    \centering
    \includegraphics[width=0.9\linewidth]{Zawadzcy_podział_ziemi_2.JPG}
    \captionsetup{format=hang}
    \caption{Fragment wpisu w~księdze hipotecznej wsi \emph{Desna} informujący
    o~ostatecznym podziale majątku Marcina i~Marianny Zawadzkich - część 2.
    \cite{desna1}.}
    \label{fig:mzawadzki_1944_2}
\end{figure}

\begin{figure}[!ht]
    \vspace*{0.5cm}
    \centering
    \includegraphics[width=1.0\linewidth]{Zawadzcy_podział_ziemi_3.JPG}
    \captionsetup{format=hang}
    \caption{Fragment wpisu w~księdze hipotecznej wsi \emph{Desna} informujący
    o~ostatecznym podziale majątku Marcina i~Marianny Zawadzkich - część 3.
    \cite{desna1}.}
    \label{fig:mzawadzki_1944_3}
\end{figure}

Z~treści przywołanego aktu hipotecznego wynika również, iż~większość dzieci
Marcina i~Marianny Zawadzkich zdecydowała się na~sprzedaż swoich udziałów
w~majątku rodziców jeszcze przed formalnym zamknięciem postępowania spadkowego:

\begin{itemize}
\item Józef\indeksosoby{Józef}{Zawadzki}{1882}{Anielina} oraz Stanisława
\indeksosoby[Zawadzka]{Stanisława}{Wróbel}{1885}{Anielina} sprzedali swoje
części ojcowizny ich młodszemu bratu
Feliksowi\indeksosoby{Feliks}{Zawadzki}{1900}{Desno}, który zapłacił
im~po~1000 ówczesnych złotych,
\item Cecylia\indeksosoby{Cecylia}{Zawadzka}{1908}{Desno} sprzedała 1/9
majątku rodziców swojemu starszemu bratu Pawłowi za~kwotę 1400 złotych,
\item Wincenty\indeksosoby{Wincenty}{Zawadzki}{1887}{Anielina}, 
Jan\indeksosoby{Jan}{Zawadzki}{1889}{Anielina}
oraz~Antoni\indeksosoby{Antoni}{Zawadzki}{1899}{Desno} wszystkie swoje prawa
do~spadku po~rodzicach sprzedali
Stanisławowi\indeksosoby{Stanisław}{Gomulski}{1892}{Mikanów}
i~Mariannie\indeksosoby[Kowalczyk]{Marianna}{Gomulska}{1897}{Lubomin} (z~domu
Kowalczyk) Gomulskim\footnote{Stanisław był szóstym synem Michała i~Marianny
Gomulskich współzałożycieli spółki \emph{Desna}.}, za~kwotę 3900 złotych.
\end{itemize}

W~konsekwencji jedna trzecia majątku małżonków Zawadzkich przypadła ich ósmemu
synowi Feliksowi (około 2,7~ha), dwie dziewiąte piątemu synowi Pawłowi (około
1,9~ha) oraz jedna dziewiąta drugiej córce
Magdalenie\indeksosoby[Zawadzka][Gomulska]{Magdalena}{Poziemska}{1895}{Cisie}
(około 0,9~ha). Pozostała część majątku Zawadzkich na~Deśnie trafiła
do~rodziny Gomulskich~-- do Stanisława i~jego żony Marianny (patrz: ryc.
\ref{fig:mzawadzki_1944_3}).

\newpage
\begin{figure}[!ht]
\centering
\begin{tikzpicture}[
  % Definiujemy stałe odstępy
  node distance=1cm and 1cm, 
  every node/.style={draw, rectangle, align=center, minimum width=3.8cm,
  inner sep=5pt}]

% --- POKOLENIE 0 ---
\node (Osiowa0) [draw=none, minimum width=0pt, inner sep=0pt] at (0,0) {};
\node (Jan) [left=1.5cm of Osiowa0] {Jan Zawadzki \\ {
    \footnotesize \emph{?, ? r.}}};
\node (Apolonia) [right=1.5cm of Osiowa0] {Apolonia Bąk \\ {
    \footnotesize \emph{?, ? r.}}};

\draw[-] (Jan) -- (Apolonia)
    node[midway, below, draw=none, minimum width=0pt,
    font=\scriptsize] (Slub0) {\emph{?, ? r.}};

% --- POKOLENIE 1 ---
\node (Osiowa1) [draw=none, minimum width=0pt, inner sep=0pt,
below=1.5cm of Slub0] {};
\node (Jozef) [left=1.5cm of Osiowa1] {Józef Zawadzki \\ {
    \footnotesize \emph{Zakole, 1815 r.}}};
\node (Marianna) [right=1.5cm of Osiowa1] {Marianna Górka \\ {
    \footnotesize \emph{Rudno, 1827 r.}}};

\draw[-] (Jozef) -- (Marianna)
    node[midway, above, draw=none, minimum width=0pt,
    font=\scriptsize] (Slub1_Gora) { }
    node[midway, below, draw=none, minimum width=0pt,
    font=\scriptsize] (Slub1) {\emph{Mińsk, 1855 r.}};

\draw[->] (Slub0) -- (Slub1_Gora);

% --- POKOLENIE 2 ---
\node (Osiowa2) [draw=none, minimum width=0pt, inner sep=0pt,
below=1.5cm of Slub1] {};
\node (Marcin) [left=1.5cm of Osiowa2] {Marcin Zawadzki \\ {
    \footnotesize \emph{Anielina, 1856 r.}}};
\node (Marianna2) [right=1.5cm of Osiowa2] {Marianna Zawadzka \\ {
    \footnotesize \emph{Targówka, 1864 r.}}};
    
\draw[-] (Marcin) -- (Marianna2)
    node[midway, above, draw=none, minimum width=0pt,
    font=\scriptsize] (Slub2_Gora) { }
    node[midway, below, draw=none, minimum width=0pt,
    font=\scriptsize] (Slub2) {\emph{Mińsk, 1881 r.}};

\draw[->] (Slub1) -- (Slub2_Gora);

% --- POKOLENIE 3 ---
\node (Osiowa3) [draw=none, minimum width=0pt, inner sep=0pt,
below=1.5cm of  Slub2] {};
\node (Magdalena) [left=1.5cm of Osiowa3] {Magdalena Zawadzka \\ {
    \footnotesize \emph{Cisie, 1895 r.}}};
\node (Michal) [right=1.5cm of Osiowa3] {Michał Poziemski \\ {
    \footnotesize \emph{Chobot, 1901 r.}}};

\draw[-] (Magdalena) -- (Michal)
    node[midway, above, draw=none, minimum width=0pt,
    font=\scriptsize] (Slub3_Gora) { }
    node[midway, below, draw=none, minimum width=0pt,
    font=\scriptsize] (Slub3) {\emph{Długa Kościelna, 1929 r.}};

\draw[->] (Slub2) -- (Slub3_Gora);

% --- POKOLENIE 4 ---
\node (Jadwiga) [below=1.5cm of Slub3] {Jadwiga Leokadia Poziemska \\
{\footnotesize \emph{Desno, 1935 r.}}};

\draw[->] (Slub3) -- (Jadwiga);

\end{tikzpicture}
\end{figure}

Spośród czternastki dzieci małżonków Zawadzkich tylko potomkowie ich córki
Magdaleny żyją dalej na~Deśnie~-- jednym z~nich jest autor książki. Pozostałe
dzieci albo zmarły we wczesnym dzieciństwie albo wyprowadziły się z~rodzinnej
miejscowości, której ich rodzice byli współzałożycielami.